\documentclass[10pt,a4paper]{amsart}
\usepackage{amsmath,amssymb,amsthm,graphicx}
% Retained preamble for the Stacks Project chapter morphisms.tex (Morphisms of Schemes).
% The chapter file pinned-original-source.tex carries only \input{preamble} on its first line, so no
% preamble text lives inside the chapter itself.  The Stacks Project uses one single shared preamble for
% every chapter; the snapshot stages it verbatim in sources/stacks-chow/preamble.tex.  The block below is
% copied mechanically, line for line, from that file's lines 14-18 (the Xy-pic setup) and lines 154-217
% (the native theorem declarations together with the mathematical macros).  Nothing was added, removed or
% redefined.  The wrapper class and the amsmath/amssymb/amsthm/graphicx base packages are supplied by the
% assembler, so the class/\externaldocument part of the shared preamble is not repeated here.
% For commutative diagrams we use Xy-pic
\usepackage[all]{xy}

% We use 2cell for 2-commutative diagrams.
\xyoption{2cell}
% Theorem environments.
%
\theoremstyle{plain}
\newtheorem{theorem}[subsection]{Theorem}
\newtheorem{proposition}[subsection]{Proposition}
\newtheorem{lemma}[subsection]{Lemma}

\theoremstyle{definition}
\newtheorem{definition}[subsection]{Definition}
\newtheorem{example}[subsection]{Example}
\newtheorem{exercise}[subsection]{Exercise}
\newtheorem{situation}[subsection]{Situation}

\theoremstyle{remark}
\newtheorem{remark}[subsection]{Remark}
\newtheorem{remarks}[subsection]{Remarks}

\numberwithin{equation}{subsection}

% Macros
%
\def\lim{\mathop{\mathrm{lim}}\nolimits}
\def\colim{\mathop{\mathrm{colim}}\nolimits}
\def\Spec{\mathop{\mathrm{Spec}}}
\def\Hom{\mathop{\mathrm{Hom}}\nolimits}
\def\Ext{\mathop{\mathrm{Ext}}\nolimits}
\def\SheafHom{\mathop{\mathcal{H}\!\mathit{om}}\nolimits}
\def\SheafExt{\mathop{\mathcal{E}\!\mathit{xt}}\nolimits}
\def\Sch{\mathit{Sch}}
\def\Mor{\mathop{\mathrm{Mor}}\nolimits}
\def\Ob{\mathop{\mathrm{Ob}}\nolimits}
\def\Sh{\mathop{\mathit{Sh}}\nolimits}
\def\NL{\mathop{N\!L}\nolimits}
\def\CH{\mathop{\mathrm{CH}}\nolimits}
\def\proetale{{pro\text{-}\acute{e}tale}}
\def\etale{{\acute{e}tale}}
\def\QCoh{\mathit{QCoh}}
\def\Ker{\mathop{\mathrm{Ker}}}
\def\Im{\mathop{\mathrm{Im}}}
\def\Coker{\mathop{\mathrm{Coker}}}
\def\Coim{\mathop{\mathrm{Coim}}}

% Boxtimes
%
\DeclareMathSymbol{\boxtimes}{\mathbin}{AMSa}{"02}

%
% Macros for moduli stacks/spaces
%
\def\QCohstack{\mathcal{QC}\!\mathit{oh}}
\def\Cohstack{\mathcal{C}\!\mathit{oh}}
\def\Spacesstack{\mathcal{S}\!\mathit{paces}}
\def\Quotfunctor{\mathrm{Quot}}
\def\Hilbfunctor{\mathrm{Hilb}}
\def\Curvesstack{\mathcal{C}\!\mathit{urves}}
\def\Polarizedstack{\mathcal{P}\!\mathit{olarized}}
\def\Complexesstack{\mathcal{C}\!\mathit{omplexes}}
% \Pic is the operator that assigns to X its picard group, usage \Pic(X)
% \Picardstack_{X/B} denotes the Picard stack of X over B
% \Picardfunctor_{X/B} denotes the Picard functor of X over B
\def\Pic{\mathop{\mathrm{Pic}}\nolimits}
\def\Picardstack{\mathcal{P}\!\mathit{ic}}
\def\Picardfunctor{\mathrm{Pic}}
\def\Deformationcategory{\mathcal{D}\!\mathit{ef}}

% --- appended 2026-10-12 (Chief): the project own definitions of the reference/slogan/history
% display environments, copied mechanically from the shared Stacks preamble lines 8-12
% (\usepackage{verbatim} plus the three \newenvironment lines, which the chapter preamble keeps
% verbatim from the project preamble). Nothing added, removed or redefined: these environments
% carry no mathematical content and the project itself typesets them as comments. Without these
% three lines a counted statement whose own text uses them fails with Environment slogan undefined.
\usepackage{verbatim}
\newenvironment{reference}{\comment}{\endcomment}
\newenvironment{slogan}{\comment}{\endcomment}
\newenvironment{history}{\comment}{\endcomment}
\usepackage{cleveref}
\crefname{originallemma}{original lemma}{original lemmas}
\Crefname{originallemma}{Original lemma}{Original lemmas}
\makeatletter
\newlabel{algebra-lemma-going-up-closed}{{00HZ}{source}{}{sourcepointer-1}{}}
\newlabel{algebra-lemma-going-up-closed@cref}{{[originallemma][0][]00HZ}{[1][0][]source}{}{}{}}
\makeatother
\crefname{originallemma}{original lemma}{original lemmas}
\Crefname{originallemma}{Original lemma}{Original lemmas}
\makeatletter
\newlabel{algebra-lemma-integral-going-up}{{00GU}{source}{}{sourcepointer-2}{}}
\newlabel{algebra-lemma-integral-going-up@cref}{{[originallemma][0][]00GU}{[1][0][]source}{}{}{}}
\makeatother
\crefname{originalsection}{original section}{original sections}
\Crefname{originalsection}{Original section}{Original sections}
\makeatletter
\newlabel{example-scheme-theoretic-image}{{056A}{source}{}{sourcepointer-3}{}}
\newlabel{example-scheme-theoretic-image@cref}{{[originalsection][0][]056A}{[1][0][]source}{}{}{}}
\makeatother
\crefname{originallemma}{original lemma}{original lemmas}
\Crefname{originallemma}{Original lemma}{Original lemmas}
\makeatletter
\newlabel{lemma-quasi-compact-scheme-theoretic-image}{{01R8}{source}{}{sourcepointer-4}{}}
\newlabel{lemma-quasi-compact-scheme-theoretic-image@cref}{{[originallemma][0][]01R8}{[1][0][]source}{}{}{}}
\makeatother
\title{A morphism is integral exactly when it is affine and universally closed}
\author{The Stacks Project Authors}
\date{stacks-project:morphisms.tex:local-snapshot-20261012; GNU FDL 1.2 or later}
\begin{document}
\maketitle
\noindent\textit{Editorial note.} All delivered mathematical text is the author's own, copied from the pinned Stacks Project chapter file pinned-original-source.tex; no statement or proof step was rewritten, reordered or removed.\par\medskip
\noindent\textit{Editorial note.} Following the multi-claim convention for this source, source\_url is the source-level URL https://stacks.math.columbia.edu/ and the source label lemma-integral-universally-closed is carried in claim\_key and here rather than in a per-claim tag URL.\par\medskip
\noindent\textit{Editorial note.} The preamble is the chapter's shared Stacks Project preamble, retained as sources/stacks-morphisms/retained-preamble.tex: a line-for-line excerpt of the shared Stacks preamble staged with this snapshot. No macro was invented or redefined, and this card adds no preamble text of its own.\par\medskip
\noindent\textit{Editorial note.} Statement span 11053-11064 and proof span 11066-11118 come from the halfA pool card operations/pool-cards/stacks-morphisms-halfA.json and were re-read in the pinned chapter file before this card was written.\par\medskip
\noindent\textit{Editorial note.} Re-checked by hand before carding: the lane's thin-proof prefilter flagged this candidate, and reading the original shows the hit comes only from short incidental phrases ('affine by definition' in the first paragraph) while the second half is a full coefficient computation plus a base change and pullback-square argument. Per the fleet rule of 2026-10-14 the hit on a proof of more than 800 bytes was re-read by hand and the reason is recorded here.\par\medskip
\noindent\textit{Editorial note.} The statement carries an author slogan ('integral equals affine plus universally closed.'); the shared preamble defines the slogan environment as a comment environment, so it is not typeset, and it is preserved verbatim in the counted span.\par\medskip
\noindent\textit{Editorial note.} 4 cross-references occur inside the counted spans; they are listed in external\_references as navigation pointers only.\par\medskip
\noindent\textit{Allowed standard prerequisites (editorial).}\begin{itemize}
\item An integral ring map satisfies going up, and a morphism satisfying going up between affine schemes is closed, cited from the algebra chapter. Cited at proof lines 11066-11070. Gives closedness of an integral morphism and hence, with base change, universal closedness.
\item The description of the scheme theoretic image of the closed subscheme cut out by a relation as the closed subscheme defined by the kernel of the corresponding ring map, and the resulting surjectivity, cited from the same chapter and an example of the schemes chapter. Cited at proof lines 11098-11102. Identifies the image of the locally closed piece and shows it surjects onto the closed subscheme defined by the kernel.
\item The identification of a pullback in the displayed square as trivial, because one map sends the variable to an invertible element while the other sends it to zero. Used at proof lines 11112-11116. Makes the bottom left corner of the diagram empty, which forces the bottom middle to be empty by surjectivity.
\end{itemize}\medskip
\ifdefined\hypertarget\hypertarget{sourcepointer-1}{}\fi
\noindent\textit{Source reference.} Original lemma 00HZ: focused-seeds/AGP/sources/stacks-snapshot-20261011/tags.txt, source lines 661--661. Same-chapter reference used in the retained author text at source line(s) 11071. It is supplied only as a navigation pointer so that the retained text compiles; the counted item delivers the author's own statement and the author's complete proof, and no part of the referenced item is copied in or substituted for the counted argument.\par\medskip
\ifdefined\hypertarget\hypertarget{sourcepointer-2}{}\fi
\noindent\textit{Source reference.} Original lemma 00GU: focused-seeds/AGP/sources/stacks-snapshot-20261011/tags.txt, source lines 620--620. Same-chapter reference used in the retained author text at source line(s) 11070, 11143. It is supplied only as a navigation pointer so that the retained text compiles; the counted item delivers the author's own statement and the author's complete proof, and no part of the referenced item is copied in or substituted for the counted argument.\par\medskip
\ifdefined\hypertarget\hypertarget{sourcepointer-3}{}\fi
\noindent\textit{Source reference.} Original section 056A: focused-seeds/AGP/sources/stacks-snapshot-20261011/tags.txt, source lines 6703--6703. Same-chapter reference used in the retained author text at source line(s) 11099. It is supplied only as a navigation pointer so that the retained text compiles; the counted item delivers the author's own statement and the author's complete proof, and no part of the referenced item is copied in or substituted for the counted argument.\par\medskip
\ifdefined\hypertarget\hypertarget{sourcepointer-4}{}\fi
\noindent\textit{Source reference.} Original lemma 01R8: focused-seeds/AGP/sources/stacks-snapshot-20261011/tags.txt, source lines 2293--2293. Same-chapter reference used in the retained author text at source line(s) 891, 903, 920, 928, 1007, 1084, 1167, 1169, 1369, 4982, 10102, 11100, 14144. It is supplied only as a navigation pointer so that the retained text compiles; the counted item delivers the author's own statement and the author's complete proof, and no part of the referenced item is copied in or substituted for the counted argument.\par\medskip
\begin{lemma}
\label{lemma-integral-universally-closed}
\begin{slogan}
integral $=$ affine $+$ universally closed
\end{slogan}
Let $f : X \to S$ be a morphism of schemes.
The following are equivalent
\begin{enumerate}
\item $f$ is integral, and
\item $f$ is affine and universally closed.
\end{enumerate}
\end{lemma}

\begin{proof}
Assume (1). An integral morphism is affine by definition.
A base change of an integral morphism is integral so in order to prove (2)
it suffices to show that an integral morphism is closed.
This follows from Algebra, Lemmas \ref{algebra-lemma-integral-going-up}
and \ref{algebra-lemma-going-up-closed}.

\medskip\noindent
Assume (2). We may assume $f$ is the morphism
$f : \Spec(A) \to \Spec(R)$ coming from a ring map
$R \to A$. Let $a$ be an element of $A$. We have to show
that $a$ is integral over $R$, i.e. that in the kernel
$I$ of the map $R[x] \to A$ sending $x$ to $a$ there is
a monic polynomial. Consider the ring
$B = A[x]/(ax -1)$ and let $J$ be the kernel of the
composition $R[x]\to A[x] \to B$. If $f\in J$ there exists
$q\in A[x]$ such that $f = (ax-1)q$ in $A[x]$ so if
$f = \sum_i f_ix^i$ and $q = \sum_iq_ix^i$, for all $i \geq 0$
we have $f_i = aq_{i-1} - q_i$. For $n \geq \deg q + 1$ the polynomial
$$
\sum\nolimits_{i \geq 0} f_i x^{n - i} =
\sum\nolimits_{i \geq 0} (a q_{i - 1} - q_i) x^{n - i} =
(a - x) \sum\nolimits_{i \geq 0} q_i x^{n - i - 1}
$$
is clearly in $I$; if $f_0 = 1$ this polynomial is also monic, so we
are reduced to prove that $J$ contains a polynomial with constant term $1$.
We do it by proving $\Spec(R[x]/(J + (x))$ is empty.


\medskip\noindent
Since $f$ is universally closed the base change $\Spec(A[x]) \to \Spec(R[x])$
is closed. Hence the image of the closed subset $\Spec(B) \subset \Spec(A[x])$
is the closed subset $\Spec(R[x]/J) \subset \Spec(R[x])$, see
Example \ref{example-scheme-theoretic-image} and
Lemma \ref{lemma-quasi-compact-scheme-theoretic-image}.
In particular $\Spec(B) \to \Spec(R[x]/J)$ is surjective. Consider the
following diagram where every square is a pullback:
$$
\xymatrix{
\Spec(B) \ar@{->>}[r]^g &
\Spec(R[x]/J)  \ar[r] &
\Spec(R[x])\\
\emptyset \ar[u] \ar[r] &
\Spec(R[x]/(J + (x)))\ar[u] \ar[r] &
\Spec(R) \ar[u]^0
}
$$
The bottom left corner is empty because it is the spectrum of
$R\otimes_{R[x]} B$ where the map $R[x]\to B$ sends $x$ to an
invertible element and $R[x]\to R$ sends $x$ to $0$. Since $g$
is surjective this implies $\Spec(R[x]/(J + (x)))$ is empty, as
we wanted to show.
\end{proof}

\end{document}
